\chapter{2008年四川卷（文）}

\section{单选题}

\begin{problem}
设集合 \(U = \{1, 2, 3, 4, 5\}\)，\(A = \{1, 2, 3\}\)，\(B = \{2, 3, 4\}\)，则 \(\complement_U(A \cap B) =\)（\quad）

\choices
    {\(\{2, 3\}\)}
    {\(\{1, 4, 5\}\)}
    {\(\{4, 5\}\)}
    {\(\{1, 5\}\)}
\end{problem}

\begin{answer}
B
\end{answer}

\begin{solution}
由 \(A\cap B=\{2,3\}\)，直接取全集 \(U\) 中除去 \(2\)，\(3\) 之外的元素，得 \(\complement_U(A\cap B)=\{1,4,5\}\). 故选 B.
\end{solution}

\begin{problem}
函数 \(y = \ln(2x + 1)\)（\(x > -\frac{1}{2}\)）的反函数是（\quad）

\choices
    {\(y = \frac{1}{2}\e^{x} - 1\)（\(x \in \R\)）}
    {\(y = \e^{2x} - 1\)（\(x \in \R\)）}
    {\(y = \frac{1}{2}(\e^{x} - 1)\)（\(x \in \R\)）}
    {\(y = \e^{\frac{x}{2}} - 1\)（\(x \in \R\)）}
\end{problem}

\begin{answer}
C
\end{answer}

\begin{solution}
设原函数的反函数为 \(y=g(x)\). 由 \(y=\ln(2x+1)\)，得 \(\e^{y}=2x+1\)，从而 \(x=\frac{\e^{y}-1}{2}\). 交换 \(x\)，\(y\) 后，反函数为 \(y=\frac{1}{2}(\e^{x}-1)\). 原函数的值域为 \(\R\)，所以反函数的定义域为 \(x\in\R\). 故选 C.
\end{solution}

\begin{problem}
设平面向量 \(\bs{a} = (3, 5)\)，\(\bs{b} = (-2, 1)\)，则 \(\bs{a} - 2\bs{b} =\)（\quad）

\choices
    {\((7, 3)\)}
    {\((7, 7)\)}
    {\((1, 7)\)}
    {\((1, 3)\)}
\end{problem}

\begin{answer}
A
\end{answer}

\begin{solution}
\(\bs{a}-2\bs{b}=\left(3,5\right)-2\left(-2,1\right)=\left(3,5\right)-\left(-4,2\right)=\left(7,3\right)\). 故选 A.
\end{solution}

\begin{problem}
\((\tan x + \cot x)\cos^{2}x =\)（\quad）

\choices
    {\(\tan x\)}
    {\(\sin x\)}
    {\(\cos x\)}
    {\(\cot x\)}
\end{problem}

\begin{answer}
D
\end{answer}

\begin{solution}
在表达式有意义的范围内，\(\sin x\ne0\)，\(\cos x\ne0\). 因而
\((\tan x+\cot x)\cos^{2}x=\left(\frac{\sin x}{\cos x}+\frac{\cos x}{\sin x}\right)\cos^{2}x=\frac{\sin^{2}x+\cos^{2}x}{\sin x}\cos x=\cot x\). 故选 D.
\end{solution}

\begin{problem}
不等式 \(|x^{2} - x| < 2\) 的解集为（\quad）

\choices
    {\((-1, 2)\)}
    {\((-1, 1)\)}
    {\((-2, 1)\)}
    {\((-2, 2)\)}
\end{problem}

\begin{answer}
A
\end{answer}

\begin{solution}
由 \(\left\lvert x^{2}-x\right\rvert<2\)，得 \(-2<x^{2}-x<2\). 其中
\(x^{2}-x>-2\) 等价于 \(x^{2}-x+2=\left(x-\frac12\right)^{2}+\frac74>0\)，恒成立；而 \(x^{2}-x<2\) 等价于 \(x^{2}-x-2<0\)，即 \((x-2)(x+1)<0\)，所以 \(-1<x<2\). 故选 A.
\end{solution}

\begin{problem}
直线 \(y = 3x\) 绕原点逆时针旋转 \(90^{\circ}\)，再向右平移 \(1\) 个单位，所得到的直线为（\quad）

\choices
    {\(y = -\frac{1}{3}x + \frac{1}{3}\)}
    {\(y = -\frac{1}{3}x + 1\)}
    {\(y = 3x - 3\)}
    {\(y = 3x + 1\)}
\end{problem}

\begin{answer}
A
\end{answer}

\begin{solution}
取原直线上一点 \(\left(t,3t\right)\). 绕原点逆时针旋转 \(90^{\circ}\) 后变为 \(\left(-3t,t\right)\)，故旋转后的直线满足 \(Y=-\frac13X\). 再向右平移 \(1\) 个单位，令 \(x=X+1\)，\(y=Y\)，得 \(y=-\frac13(x-1)=-\frac13x+\frac13\). 故选 A.
\end{solution}

\begin{problem}
\(\triangle ABC\) 的三个内角 \(A\)，\(B\)，\(C\) 的对边边长分别是 \(a\)，\(b\)，\(c\). 若 \(a = \frac{\sqrt{5}}{2}b\)，\(A = 2B\)，则 \(\cos B =\)（\quad）

\choices
    {\(\frac{\sqrt{5}}{3}\)}
    {\(\frac{\sqrt{5}}{4}\)}
    {\(\frac{\sqrt{5}}{5}\)}
    {\(\frac{\sqrt{5}}{6}\)}
\end{problem}

\begin{answer}
B
\end{answer}

\begin{solution}
由正弦定理，\(\frac{a}{b}=\frac{\sin A}{\sin B}=\frac{\sin 2B}{\sin B}=2\cos B\). 结合 \(a=\frac{\sqrt5}{2}b\)，得 \(2\cos B=\frac{\sqrt5}{2}\)，所以 \(\cos B=\frac{\sqrt5}{4}\). 故选 B.
\end{solution}

\begin{problem}
设 \(M\) 是球 \(O\) 的半径 \(OP\) 的中点，分别过 \(M\)，\(O\) 作垂直于 \(OP\) 的平面，截球面得到两个圆，则这两个圆的面积比值为（\quad）

\choices
    {\(\frac{1}{4}\)}
    {\(\frac{1}{2}\)}
    {\(\frac{2}{3}\)}
    {\(\frac{3}{4}\)}
\end{problem}

\begin{answer}
D
\end{answer}

\begin{solution}
设球的半径为 \(R\). 过球心 \(O\) 的截面圆半径为 \(R\)，面积为 \(S_O=\pi R^{2}\). 由于 \(M\) 是 \(OP\) 的中点，\(OM=\frac{R}{2}\)，过 \(M\) 的截面圆半径 \(r\) 满足 \(r^{2}=R^{2}-OM^{2}=R^{2}-\frac{R^{2}}{4}=\frac34R^{2}\). 因而 \(S_M=\frac34\pi R^{2}\)，所以两个圆的面积比为 \(\frac{S_M}{S_O}=\frac34\). 故选 D.
\end{solution}

\begin{problem}
设定义在 \(\R\) 上的函数 \(f(x)\) 满足 \(f(x) \cdot f(x + 2) = 13\)，若 \(f(1) = 2\)，则 \(f(99) =\)（\quad）

\choices
    {\(13\)}
    {\(2\)}
    {\(\frac{13}{2}\)}
    {\(\frac{2}{13}\)}
\end{problem}

\begin{answer}
C
\end{answer}

\begin{solution}
由 \(f(x)f(x+2)=13\)，知 \(f(x)\ne0\)，且 \(f(x+2)=\frac{13}{f(x)}\). 再将 \(x\) 换成 \(x+2\)，得 \(f(x+4)=\frac{13}{f(x+2)}=f(x)\)，所以 \(f\) 以 \(4\) 为一个周期. 因为 \(99=1+4\times24+2\)，故 \(f(99)=f(3)=\frac{13}{f(1)}=\frac{13}{2}\). 故选 C.
\end{solution}

\begin{problem}
设直线 \(l \subset\) 平面 \(\alpha\)，过平面 \(\alpha\) 外一点 \(A\) 与 \(l\)，\(\alpha\) 都成 \(30^{\circ}\) 角的直线有且只有（\quad）

\choices
    {\(1\) 条}
    {\(2\) 条}
    {\(3\) 条}
    {\(4\) 条}
\end{problem}

\begin{answer}
B
\end{answer}

\begin{solution}
取单位方向向量 \(\bs{u}=\left(u_{1},u_{2},u_{3}\right)\)，其中 \(u_{3}\) 沿平面 \(\alpha\) 的法向，且取 \(l\) 的方向为第一个坐标轴. 直线与平面 \(\alpha\) 成 \(30^{\circ}\) 角，故 \(\lvert u_{3}\rvert=\sin30^{\circ}=\frac12\)；直线与 \(l\) 成 \(30^{\circ}\) 角，故 \(\lvert u_{1}\rvert=\cos30^{\circ}=\frac{\sqrt3}{2}\). 由 \(u_{1}^{2}+u_{2}^{2}+u_{3}^{2}=1\)，得 \(u_{2}=0\). 取 \(u_{1}>0\) 消去方向相反的重复表示后，\(u_{3}=\frac12\) 或 \(u_{3}=-\frac12\)，对应两条不同的直线. 故选 B.
\end{solution}

\begin{problem}
已知双曲线 \(C: \frac{x^{2}}{9} - \frac{y^{2}}{16} = 1\) 的左，右焦点分别为 \(F_{1}\)，\(F_{2}\)，\(P\) 为 \(C\) 的右支上一点，且 \(|PF_{2}| = |F_{1}F_{2}|\)，则 \(\triangle PF_{1}F_{2}\) 的面积等于（\quad）

\choices
    {\(24\)}
    {\(36\)}
    {\(48\)}
    {\(96\)}
\end{problem}

\begin{answer}
C
\end{answer}

\begin{solution}
双曲线的半长轴、半短轴分别为 \(a=3\)，\(b=4\)，故半焦距 \(c=\sqrt{a^{2}+b^{2}}=5\)，从而 \(F_{1}F_{2}=2c=10\). \(P\) 在右支上，满足 \(PF_{1}-PF_{2}=2a=6\)；结合 \(PF_{2}=10\)，得 \(PF_{1}=16\). 于是三角形 \(PF_{1}F_{2}\) 的三边为 \(16\)，\(10\)，\(10\)，半周长为 \(18\). 由海伦公式，面积为 \(\sqrt{18\times2\times8\times8}=48\). 故选 C.
\end{solution}

\begin{problem}
若三棱柱的一个侧面是边长为 \(2\) 的正方形，另外两个侧面都是有一个内角为 \(60^{\circ}\) 的菱形，则该棱柱的体积为（\quad）

\choices
    {\(\sqrt{2}\)}
    {\(2\sqrt{2}\)}
    {\(3\sqrt{2}\)}
    {\(4\sqrt{2}\)}
\end{problem}

\begin{answer}
B
\end{answer}

\begin{solution}
设三棱柱为 \(ABC-A_{1}B_{1}C_{1}\)，令 \(\bs{v}=\overrightarrow{AA_{1}}\). 一个侧面为边长 \(2\) 的正方形，取其底边为 \(AB\)，则 \(\lvert AB\rvert=\lvert\bs{v}\rvert=2\)，且 \(\bs{v}\perp AB\). 另外两个侧面为菱形，所以 \(\lvert AC\rvert=\lvert BC\rvert=2\). 因而底面三角形 \(ABC\) 为边长 \(2\) 的正三角形，面积为 \(\sqrt3\).

在底面内取坐标 \(A=(0,0,0)\)，\(B=(2,0,0)\)，\(C=(1,\sqrt3,0)\)，设 \(\bs{v}=(0,u,w)\). 由于侧面菱形有一个内角为 \(60^{\circ}\)，\(\lvert\bs{v}\cdot\overrightarrow{AC}\rvert=\lvert\bs{v}\rvert\lvert\overrightarrow{AC}\rvert\cos60^{\circ}=2\)，所以 \(\lvert u\rvert\sqrt3=2\)，即 \(u^{2}=\frac43\). 又 \(\lvert\bs{v}\rvert^{2}=4\)，得 \(w^{2}=4-\frac43=\frac83\). 棱柱的高为 \(\lvert w\rvert=2\sqrt{\frac23}\)，故体积为 \(\sqrt3\cdot2\sqrt{\frac23}=2\sqrt2\). 故选 B.
\end{solution}

\section{填空题}

\begin{problem}
\((1 + 2x)^{3}(1 - x)^{4}\) 展开式中 \(x\) 的系数为\fillinblank{}.
\end{problem}

\begin{answer}
\(2\)
\end{answer}

\begin{solution}
要得到 \(x\) 项，只有两种取法：从 \((1+2x)^{3}\) 取一次 \(x\)，从 \((1-x)^{4}\) 取常数；或从前者取常数，从后者取一次 \(x\). 因此 \(x\) 的系数为 \(\mathrm C_{3}^{1}\cdot2+\mathrm C_{4}^{1}\cdot(-1)=6-4=2\).
\end{solution}

\begin{problem}
已知直线 \(l: x - y + 4 = 0\) 与圆 \(C: (x - 1)^{2} + (y - 1)^{2} = 2\)，则 \(C\) 上各点到 \(l\) 的距离的最小值为\fillinblank{}.
\end{problem}

\begin{answer}
\(\sqrt{2}\)
\end{answer}

\begin{solution}
圆心为 \(O=\left(1,1\right)\)，半径为 \(r=\sqrt2\). 圆心到直线 \(l\) 的距离为 \(d=\frac{\lvert1-1+4\rvert}{\sqrt{1^{2}+(-1)^{2}}}=2\sqrt2\). 因为 \(d>r\)，圆上各点到直线的最小距离等于 \(d-r=2\sqrt2-\sqrt2=\sqrt2\).
\end{solution}

\begin{problem}
从甲，乙等 \(10\) 名同学中挑选 \(4\) 名参加某项公益活动，要求甲，乙中至少有 \(1\) 人参加，则不同的挑选方法有\fillinblank{}种.
\end{problem}

\begin{answer}
\(140\)
\end{answer}

\begin{solution}
从 \(10\) 名同学中任选四人的方法数为 \(\mathrm C_{10}^{4}\). 不满足条件的情形是甲、乙都不参加，此时从其余 \(8\) 人中选 \(4\) 人，有 \(\mathrm C_{8}^{4}\) 种. 所求方法数为 \(\mathrm C_{10}^{4}-\mathrm C_{8}^{4}=210-70=140\) 种.
\end{solution}

\begin{problem}
设数列 \(\{a_{n}\}\) 中，\(a_{1} = 2\)，\(a_{n+1} = a_{n} + n + 1\)，则通项 \(a_{n} =\)\fillinblank{}.
\end{problem}

\begin{answer}
\(\frac{n^{2} + n + 2}{2}\)
\end{answer}

\begin{solution}
由递推式逐项相加，得 \(a_{n}=a_{1}+\sum_{k=1}^{n-1}(k+1)=2+\frac{n(n-1)}{2}+n-1=\frac{n^{2}+n+2}{2}\).
\end{solution}

\section{解答题}

\begin{problem}
求函数 \(y = 7 - 4\sin x\cos x + 4\cos^{2}x - 4\cos^{4}x\) 的最大值与最小值.
\end{problem}

\begin{solution}
令 \(t=\sin x\cos x\). 由 \(t=\frac12\sin 2x\)，得 \(-\frac12\leqslant t\leqslant\frac12\). 又 \(4\cos^{2}x-4\cos^{4}x=4\sin^{2}x\cos^{2}x=4t^{2}\)，所以 \(y=4t^{2}-4t+7=4\left(t-\frac12\right)^{2}+6\). 当 \(t=\frac12\) 时，\(y\) 取得最小值 \(6\)；当 \(t=-\frac12\) 时，\(y\) 取得最大值 \(10\). 因此最大值为 \(10\)，最小值为 \(6\).
\end{solution}

\begin{problem}
设进入某商场的每一位顾客购买甲种商品的概率为 \(0.5\)，购买乙种商品的概率为 \(0.6\)，且购买甲种商品与购买乙种商品相互独立，各顾客之间购买商品也相互独立.

\begin{enumerate}
\item 求进入商场的 \(1\) 位顾客购买甲、乙两种商品中的一种的概率；
\item 求进入该商场的 \(3\) 位顾客中，至少有 \(2\) 位顾客既未购买甲种也未购买乙种商品的概率.
\end{enumerate}
\end{problem}

\begin{solution}
\begin{enumerate}
\item 设顾客购买甲、乙两种商品分别为事件 \(A\)，\(B\). 因为 \(A\)，\(B\) 相互独立，所以购买甲、乙两种商品中的一种的概率为 \(P(A\overline{B})+P(\overline{A}B)=0.5\times0.4+0.5\times0.6=0.5\).
    \item 一位顾客既未购买甲种商品也未购买乙种商品的概率为 \(P(\overline{A}\cap\overline{B})=0.5\times0.4=0.2\). 设进入商场的 \(3\) 位顾客中既未购买甲种商品也未购买乙种商品的人数为 \(X\)，则 \(X\sim B(3,0.2)\). 因此至少有 \(2\) 位顾客既未购买甲种商品也未购买乙种商品的概率为 \(P(X\geqslant2)=\mathrm C_3^2(0.2)^{2}(0.8)+(0.2)^{3}=0.104\).
\end{enumerate}
\end{solution}

\begin{problem}
如图，平面 \(ABEF \perp\) 平面 \(ABCD\)，四边形 \(ABEF\) 与 \(ABCD\) 都是直角梯形，\(\angle BAD = \angle FAB = 90^{\circ}\)，\(\overrightarrow{BC}=\frac{1}{2}\overrightarrow{AD}\)，\(\overrightarrow{BE}=\frac{1}{2}\overrightarrow{AF}\)，\(G\)，\(H\) 分别是 \(FA\)，\(FD\) 的中点.

\begin{enumerate}
\item 证明：四边形 \(BCHG\) 是平行四边形；
\item \(C\)，\(D\)，\(E\)，\(F\) 四点是否共面？为什么？
\item 设 \(AB = BE\). 证明：平面 \(ADE \perp\) 平面 \(CDE\).
\end{enumerate}

\bitmapfigure[max width=0.52\linewidth]{img/2008/sichuan_ordinary_liberal/q19_fig1.png}
\end{problem}

\begin{solution}
\begin{enumerate}
\item 在三角形 \(FAD\) 中，\(G\)，\(H\) 分别是 \(FA\)，\(FD\) 的中点，所以 \(GH\parallel AD\)，且 \(GH=\frac12AD\). 由题设 \(BC\parallel AD\)，且 \(BC=\frac12AD\)，从而 \(GH\parallel BC\)，\(GH=BC\). 因此四边形 \(BCHG\) 是平行四边形.
\item 由题设 \(\overrightarrow{BC}=\frac12\overrightarrow{AD}\)，\(\overrightarrow{BE}=\frac12\overrightarrow{AF}\)，有 \(\overrightarrow{CE}=\overrightarrow{CB}+\overrightarrow{BE}=\frac12\left(\overrightarrow{AF}-\overrightarrow{AD}\right)=\frac12\overrightarrow{DF}\). 因此 \(CE\parallel DF\)，所以直线 \(CE\) 与直线 \(DF\) 共面，故点 \(C\)，\(D\)，\(E\)，\(F\) 共面.
\item 由平面 \(ABEF\perp\) 平面 \(ABCD\)，且 \(AF\perp AB\)，得 \(AF\perp\) 平面 \(ABCD\). 以 \(A\) 为原点，分别以 \(AB\)，平面 \(ABCD\) 内垂直于 \(AB\) 的方向，\(AF\) 为三个坐标轴正方向，设 \(AB=a\)，\(BC=b\)，\(BE=c\)，则可取 \(A(0,0,0)\)，\(B(a,0,0)\)，\(C(a,b,0)\)，\(D(0,2b,0)\)，\(E(a,0,c)\). 由 \(AB=BE\) 得 \(a=c\). 于是平面 \(ADE\) 的一个法向量为 \(\overrightarrow{AD}\times\overrightarrow{AE}=(0,2b,0)\times(a,0,a)=(2ab,0,-2ab)\)，平面 \(CDE\) 的一个法向量为 \(\overrightarrow{DC}\times\overrightarrow{DE}=(a,-b,0)\times(a,-2b,a)=(-ab,-a^{2},-ab)\). 两法向量的数量积为 \(-2a^{2}b^{2}+2a^{2}b^{2}=0\)，所以平面 \(ADE\perp\) 平面 \(CDE\).
\end{enumerate}
\end{solution}

\begin{problem}
设 \(x = 1\) 和 \(x = 2\) 是函数 \(f(x) = x^{5} + ax^{3} + bx + 1\) 的两个极值点.

\begin{enumerate}
\item 求 \(a\) 和 \(b\) 的值；
\item 求 \(f(x)\) 的单调区间.
\end{enumerate}
\end{problem}

\begin{solution}
\begin{enumerate}
\item 因为 \(f(x)\) 可导，且 \(x=1\)，\(x=2\) 是极值点，所以 \(f'(1)=f'(2)=0\). 由 \(f'(x)=5x^{4}+3ax^{2}+b\)，得
\(\begin{cases}
5+3a+b=0,\\
80+12a+b=0.
\end{cases}\)
两式相减得 \(9a+75=0\)，所以 \(a=-\frac{25}{3}\)，代回得 \(b=20\).
\item 由 \(a=-\frac{25}{3}\)，\(b=20\)，得 \(f'(x)=5x^{4}-25x^{2}+20=5(x^{2}-1)(x^{2}-4)\). 当 \(x\in(-\infty,-2)\cup(-1,1)\cup(2,+\infty)\) 时，\(f'(x)>0\)；当 \(x\in(-2,-1)\cup(1,2)\) 时，\(f'(x)<0\). 因此 \(f(x)\) 的单调递增区间为 \((-\infty,-2)\)，\((-1,1)\)，\((2,+\infty)\)，单调递减区间为 \((-2,-1)\)，\((1,2)\).
\end{enumerate}
\end{solution}

\begin{problem}
已知数列 \(\{a_{n}\}\) 的前 \(n\) 项和 \(S_{n} = 2a_{n} - 2^{n}\).

\begin{enumerate}
\item 求 \(a_{3}\)，\(a_{4}\)；
\item 证明：数列 \(\{a_{n+1} - 2a_{n}\}\) 是等比数列；
\item 求 \(\{a_{n}\}\) 的通项公式.
\end{enumerate}
\end{problem}

\begin{solution}
\begin{enumerate}
\item 令 \(n=1\)，由 \(S_{1}=a_{1}=2a_{1}-2\)，得 \(a_{1}=2\). 令 \(n=2\)，由 \(S_{2}=a_{1}+a_{2}=2a_{2}-4\)，得 \(a_{2}=6\). 令 \(n=3\)，由 \(S_{3}=a_{1}+a_{2}+a_{3}=2a_{3}-8\)，得 \(a_{3}=16\). 令 \(n=4\)，由 \(S_{4}=a_{1}+a_{2}+a_{3}+a_{4}=2a_{4}-16\)，得 \(a_{4}=40\).
\item 对任意正整数 \(n\)，由 \(S_{n+1}-S_{n}=a_{n+1}\) 及已知条件，得 \(a_{n+1}=2a_{n+1}-2a_{n}-2^{n}\)，即 \(a_{n+1}-2a_{n}=2^{n}\). 因此数列 \(\{a_{n+1}-2a_{n}\}\) 的第 \(n\) 项为 \(2^{n}\)，相邻两项的比为 \(2\)，所以它是等比数列.
\item 由上式 \(a_{n+1}=2a_{n}+2^{n}\). 令 \(u_{n}=\frac{a_{n}}{2^{n-1}}\)，则 \(u_{1}=2\)，且 \(u_{n+1}=\frac{a_{n+1}}{2^{n}}=\frac{2a_{n}+2^{n}}{2^{n}}=u_{n}+1\). 所以 \(u_{n}=u_{1}+n-1=n+1\)，从而 \(a_{n}=(n+1)2^{n-1}\).
\end{enumerate}
\end{solution}

\begin{problem}
设椭圆 \(\frac{x^{2}}{a^{2}} + \frac{y^{2}}{b^{2}} = 1\)（\(a > b > 0\)）的左、右焦点分别为 \(F_{1}\)，\(F_{2}\)，离心率 \(e = \frac{\sqrt{2}}{2}\)，点 \(F_{2}\) 到右准线 \(l\) 的距离为 \(\sqrt{2}\).

\begin{enumerate}
\item 求 \(a\)，\(b\) 的值；
\item 设 \(M\)，\(N\) 是右准线 \(l\) 上两个动点，满足 \(\overrightarrow{F_{1}M} \cdot \overrightarrow{F_{2}M} = 0\). 证明：当 \(|MN|\) 取最小值时，\(\overrightarrow{F_{2}F_{1}} + \overrightarrow{F_{2}M} + \overrightarrow{F_{2}N} = 0\).
\end{enumerate}
\end{problem}

\begin{solution}
\begin{enumerate}
\item 设椭圆的半焦距为 \(c\). 由离心率定义，\(e=\frac{c}{a}=\frac{\sqrt{2}}{2}\)，所以 \(c=\frac{\sqrt{2}}{2}a\). 右准线方程为 \(x=\frac{a}{e}\)，故 \(F_{2}\) 到右准线的距离为 \(\frac{a}{e}-c=a\left(\frac1e-e\right)=\frac{\sqrt{2}}{2}a\). 由题设 \(\frac{\sqrt{2}}{2}a=\sqrt{2}\)，得 \(a=2\)，\(c=\sqrt{2}\). 再由 \(b^{2}=a^{2}-c^{2}=4-2=2\)，得 \(b=\sqrt{2}\).
\item 由上问，\(F_{1}(-\sqrt{2},0)\)，\(F_{2}(\sqrt{2},0)\)，右准线 \(l\) 的方程为 \(x=\frac{a}{e}=2\sqrt{2}\). 设 \(M(2\sqrt{2},y)\). 则
\(\overrightarrow{F_{1}M}=(3\sqrt{2},y)\)，\(\overrightarrow{F_{2}M}=(\sqrt{2},y)\)
\[
\overrightarrow{F_{1}M}\cdot\overrightarrow{F_{2}M}=3\sqrt{2}\cdot\sqrt{2}+y^{2}=6+y^{2}>0
\]
因此不存在满足题设垂直条件的实点 \(M\)，从而也不存在相应的点 \(N\) 及 \(|MN|\) 的最小值，题目第（2）问按原文件题面无法成立. 这里保留原题条件，不将 \(\overrightarrow{F_{2}M}\) 擅自改写为 \(\overrightarrow{F_{2}N}\).
\end{enumerate}
\end{solution}
